\section{Diagrams and linear filling}\label{sec:diagrams}

This section proves the combinatorial properties of the random complex
$P$ of Section~\ref{sec:model}.  An \emph{occurrence} means an included
tuple, together with its table type; copies of a cell in a diagram need
not be distinct occurrences.  All constants in this section may depend
on the fixed incidence template and on the fixed parameters, but not on
$n$.

\begin{theorem}\label{thm:diagram}
Fix sufficiently small positive $\alpha,\beta,\delta$ with
$\beta>7\alpha$, and put $f=2-\alpha$ and $p=3-\beta$.
As $n\longrightarrow\infty$, the probability tends to one that $P$ is
aspherical and its letter presentation satisfies a linear disk
isoperimetric inequality.
\end{theorem}

The central estimate will bound the number of faces in a bounded
piece of a normalized diagram by the number of its exposed sides. The expected
occurrence counts in Section~\ref{sec:model} have exponents
$2-\alpha$, $3-\beta$, and $h-\delta$ for fans, tops, and decorations
of length $2h$. The local incidence geometry compares these with the
values $2$, $3$, and $h$. After shared fan costs have been accounted
for, the differences $\alpha$, $\beta-7\alpha$, and $\delta$ supply the
strictly positive margin in the boundary estimate.

We first obtain a probability bound for every bounded collection of
distinct occurrences. We then normalize disk and sphere pictures so
that this bound also controls repeated copies of their cells. The
resulting estimate applies to arbitrary sets of faces, which lets us
partition a large picture by planar separators and sum the bounds on
its pieces. The same estimate proves linear filling for disks and
rules out nonzero two-cycles in the universal cover.

We retain every fan slot of a top occurrence when making a collection of
occurrences, even if a diagram uses only some of those slots.  This
convention is essential to the counting argument.

\subsection{Coincidence ranks}

Counting constraints while retaining repeated relator occurrences is
standard in random-group diagram arguments; see
Ollivier~\cite[Section~2]{Ollivier2004} and
Antoniuk--Friedgut--\L uczak~\cite[Section~3]{AntoniukFriedgutLuczak2017}.
We count the shared fan entries explicitly in the estimate below.

For a finite collection of top, fan, and decoration occurrences, its
\emph{field positions} are the following: the loose field of each top,
the $k_i$ fields of each fan occurrence of type $i$, and all fields of
each decoration occurrence.  A fan field is included only once when the
same fan occurrence occupies slots in several tops.  The
\emph{unsigned coincidence rank} $R$ is the number of positions minus
the number of distinct underlying generators represented at those
positions.  Thus a letter and its inverse have the same unsigned value.

\begin{lemma}\label{diagram:rank}
For each fixed integer $M$, with probability tending to one the
following holds simultaneously for every collection of at most $M$
distinct occurrences that contains all fan slots of its tops.  If the
collection has $s$ tops and $u$ fans, then
\begin{equation}\label{diagram:rank-bound}
 R\le (p-7f)s+fu+\sum_e(h_e-\delta),
\end{equation}
where the sum is over its decorations and $2h_e\in\{4,6\}$ is the
length of decoration $e$.
\end{lemma}

\begin{proof}
Fix the table types, the incidence of the fan slots, and a partition of
the field positions into unsigned-value classes.  There are only
finitely many such choices for fixed $M$.  Write $J$ for the number of
positions.  Assigning letters to a pattern of rank $R$ has at most
$C_M N^{J-R}$ possibilities: first choose one underlying generator per
class, then choose the signs of the positions.  The latter contributes
only a factor bounded in terms of $M$.

We count only choices whose occurrences are distinct within their
respective tables.  For such choices all the required Bernoulli trials
are independent.  Their product contributes the power of $N$ with
exponent
\[
 (p-7f-1)s+\sum_b(f-k_b)
       +\sum_e(h_e-\delta-2h_e).
\]
Here a shared fan occurrence appears once in the second sum, not once
per top using it.  Since
$J=s+\sum_b k_b+\sum_e2h_e$, the expected number of realizations of the
chosen pattern is at most
\[
 C_M N^{(p-7f)s+fu+\sum_e(h_e-\delta)-R}.
\]
If Equation~\eqref{diagram:rank-bound} fails, this exponent is negative.
There are finitely many exponents to consider for the fixed size bound,
so summing these expectations proves the assertion.
\end{proof}

\subsection{Spherical pictures and their reductions}

Lemma~\ref{diagram:rank} counts distinct occurrences, whereas a filling
diagram may use the same occurrence many times. We normalize the
diagram in two ways: the triangles incident to each top apex will be
grouped into polygons tracing simple macro cycles, and the cyclic
link at a fan apex will have no repeated radial label.
These properties will bound how often a field can appear near each
apex. For the asphericity argument, the normalization must also
preserve a prescribed lifted two-cycle.

An \emph{elementary picture} is a finite collection of oriented polygons with
orientation-reversing pairings of their sides, such that the resulting
closed oriented surfaces are spheres. Vertex links are taken to be
circles, with a separate source vertex for each link circle.
An interior face maps to a fan triangle, a triangle of a top cone, or a
decoration disk of $P$, with either orientation.  In the disk problem
there is one distinguished outside polygon whose boundary reads the
prescribed letter word; no extension of its interior to $P$ is required.
In the sphere problem there is no outside polygon.  Collapsed regions
mapping to vertices and free cancellations of edge paths have zero
face cost.  They can equivalently be recorded by thickening the paired
sides to strips and capping the boundary circles of the resulting
vertex regions.

Every cell has a specified map to $P$, including the index of each of
its sides.  Equal edge labels alone do not identify two side indices.
For a picture mapped to $\widetilde P$, its \emph{cellular chain} is the
sum of its signed lifted two-cells; the outside polygon is omitted.

\begin{lemma}\label{diagram:existence}
Every null letter word has a disk picture.  If
$c\ne0$ is a cellular two-cycle of $\widetilde P$, there is a finite
collection of spherical pictures with total cellular chain exactly
$c$.
\end{lemma}

\begin{proof}
The disk assertion is the cellular van Kampen construction, applied to
the finite cone and polygon complex.  Here is also a description that
applies to the sphere assertion.  Since $\widetilde P$ is simply
connected, the Hurewicz map
$\pi_2(\widetilde P)\longrightarrow H_2(\widetilde P;\mathbb Z)$ is an
isomorphism~\cite[Theorem~4.32]{Hatcher2002}.  There are no three-cells.
Consequently a sphere map
representing the class of $c$ has cellular two-chain equal to $c$, not
merely congruent to it modulo boundaries of three-cells.

Put this map in cellular general position.  One way to do so is to
choose an interior point in each two-cell, make its inverse image a
finite set of regular points, and take small disks about those
preimages.  Each such disk maps, with a sign, across that two-cell.
Retracting each punctured target cell to its boundary moves the
complement of these disks into the one-skeleton.  Inverse images of
interior edge points now give strips pairing the corresponding cell
sides; the remaining regions map to vertices.  Circular strips not
meeting cell sides may be cut and capped.  A vertex region with several
boundary components may likewise be replaced by a separate cap for
each component.  These operations are compressions along curves in an
oriented sphere.  They therefore produce a collection of oriented
spheres, not surfaces of positive genus.  They change none of the
signed cell disks or their lifts.  Components with no cell disks may
be discarded.  This gives the claimed spherical pictures and preserves
the chain $c$ exactly.

For a disk map, adjoin the distinguished outside polygon before the
same construction.  Keep that polygon as a single face throughout.
The construction preserves its word and gives one picture containing
it, possibly together with sphere components.  All the required
pictures are finite.
\end{proof}

At a vertex mapping to the apex of a top cone, every incident face is
a triangle of that cone.  Grouping its incident triangles gives a
polygon whose boundary is a closed walk in the macro graph of that top
occurrence.  We call it a \emph{top polygon}; its filling in $P$ is
always the indicated radial cone filling.  Its perimeter is the number
of these cone triangles.  The other interior faces remain fan triangles
and decoration disks.

The macro graph is the Fano incidence graph with one subdivided edge.
It has fifteen vertices.  A simple macro cycle has at most fifteen
sides, visits each point at most once, and visits at least three
points.  At a point it turns between two distinct \emph{ports}, namely
the radial labels of the fan there.  The ports are the incident line
vertices, except that the exceptional port at $D$ is $X_D$.

\begin{lemma}\label{diagram:normal-form}
A null letter word has a picture of minimum total top perimeter and,
subject to this, minimum number of remaining interior faces, with the
following properties.  The same assertion holds for a fixed nonzero
cellular two-cycle $c$, minimizing among finite collections of sphere
pictures whose total lifted cellular chain is exactly $c$.
\begin{enumerate}
\item Every top polygon is a simple macro cycle.
\item Adjacent top polygons do not use the same side of the same
top occurrence.  Two copies of one fan triangle cannot be adjacent
along equally indexed sides, and two copies of one decoration cannot
be adjacent along the same position of that decoration.
\item At each fan-apex instance, the cyclic list of radial labels has
no repetition.  A circuit with two distinct parallel steps is
permitted.  Each such circuit has at most five steps.  If it contains
no top turn, it has at least four fan triangles.
\end{enumerate}
The modifications establishing these properties preserve the outside
word and, in the sphere problem, the exact lifted cellular chain.
\end{lemma}

\begin{proof}
The nonempty class of finite pictures in
Lemma~\ref{diagram:existence} has a lexicographic minimum for the stated
pair of nonnegative integer costs.  In the sphere case fix the chain
$c$ once and for all in taking this minimum.

\emph{Simple macro cycles.}  Consider a top polygon whose macro walk repeats a template
vertex.  Its two boundary points at that repetition have the same
lifted value: both are joined to the same lifted top apex by the same
indexed cone spoke.  Split the polygon by a chord collapsed to that
vertex.  The two resulting macro walks have the same total perimeter
and the same total cone chain.  A free backtrack cancels two oppositely
oriented copies of a single lifted cone triangle and reduces total
perimeter by two.  Delete any resulting empty polygons.  Repeating
splitting and removal of backtracks therefore gives simple macro
cycles.  Collapsed chords may be put in
the picture form of Lemma~\ref{diagram:existence}; they carry no
two-chain.  This proves the first property.

\emph{Cancellation.}  Suppose two top polygons meet at the same
macro side of one occurrence.  The two elementary cone triangles
incident to that side are copies of the same triangle with opposite
orientations.  Their lifts agree, because the common lifted macro side
determines the lifted triangle and its apex.  Cancel the pair and
merge the remaining cone walks.  This reduces total top perimeter by
two.  It is an ordinary diagram cancellation: one uses the two face
interiors and a thin strip across the specified common side, so other
contacts between their boundaries cause no obstruction.  Any vertex
regions produced by the cancellation are capped as above.  A
same-template contact cannot occur between two sides of a single
simple top polygon, because a simple macro cycle uses each edge at
most once.  The identical cancellation of two fan triangles or two
decoration disks along equally indexed sides reduces the second cost
without increasing the first.  These cancellations preserve the
cellular chain, so they are forbidden by minimality in both problems.

\emph{Fan-apex links.}  Fix a source vertex $v$
mapping to a fan apex $b$.  A step in its link is either a fan triangle
or a corner of a top polygon.  Suppose two different outgoing spokes
at $v$ have the same radial label.  Each spoke is an edge from $v$ to a
vertex mapping to the letter-bouquet vertex.  Cut the side pairings of
these two edges and pair them crosswise, with the outgoing orientations
at $v$ agreeing.  This preserves the orientation of every face and all
face maps.  In terms of vertex-link cycles, it splits the link cycle of
$v$ into two.  At the far ends it either merges two distinct vertex
cycles or splits one vertex cycle.

For completeness, this switch preserves the genus-zero condition.
The numbers of faces and edges do not change.  When the two far
vertices were distinct, the total vertex count does not change, and
the affected surface remains connected: both switched edges are
incident to the merged far vertex.  When the far vertices were the
same, the total vertex count increases by two. There can be at most
two affected components after either switch. To see this, cut the two
corresponding edges of the connected dual graph before re-pairing
their four exposed half-edges. Every component of the cut graph
contains an exposed half-edge: otherwise it was already disconnected
before the cut. After re-pairing, every component therefore contains
one of the two switched edges. This argument includes loops and
parallel edges in the dual graph.
In the first case the connected oriented surface consequently has
Euler characteristic two.  In the second case the total Euler
characteristic is four, which forces exactly two oriented spheres.
All unaffected spheres are retained, and the outside polygon remains
one unchanged face.

The switch is also valid for the fixed lifted chain.  The repeated
radial label and the lifted initial fan vertex determine the same
lifted radial edge and far endpoint.  Thus the new vertex
identifications agree in $\widetilde P$, and every lifted face is
unchanged.  Each switch increases the number of fan-apex instances
while keeping the finite set of polygon corners fixed.  Repetition
therefore terminates.  The two minimized costs are unchanged, so a
newly created forbidden adjacency would still contradict their
minimality by the cancellation just described.

The resulting link uses each radial label at most once.  There are at
most five labels in any fan.  A two-step link made only of fan
triangles is the reverse traversal of one edge of the simple fan graph;
the two triangles would cancel along an equally indexed radial side.
Every other nonempty triangle-only link is a circuit in the simple
bipartite fan graph and has length at least four.  This proves the last
property.
\end{proof}

The outside word uses only letter edges, so an outside polygon never
meets a fan-apex instance.  In particular each full link considered in
Lemma~\ref{diagram:normal-form} is a cyclic link of interior faces.

\subsection{Weights on a bounded fragment}

Fix a normal picture.  A \emph{fragment} is any selection of its
interior faces; the selection need not be connected or simply
connected.  Keep a pairing of sides precisely when both incident faces
are selected.  All other selected sides are fragment boundary sides.
Let $T,Q,E$ be the numbers of selected top polygons, fan triangles, and
decoration disks, respectively, and put $A=T+Q+E$.

For a top occurrence $a$ let $t_a$ be its number of selected top
polygons, and for a decoration occurrence $e$ let $t_e$ be its selected
multiplicity.  At a fan-apex instance $z$ represented in the fragment,
let $d_z$ count selected top turns and $q_z$ count selected fan
triangles.  It is \emph{partial} if not every step of its full cyclic
link is selected.

For a fan occurrence $b$, two multiplicities must be distinguished.
The number of represented apex instances bounds the actual use of
each fan field. We must also retain this fan whenever we retain a top
of sufficiently large multiplicity that requires it. The fan weight
must therefore dominate the largest such top multiplicity. Write
\[
 r_b=\#\{z:z\text{ is represented over }b\},\qquad
 m_b=\max_{a\to b}t_a,\qquad
 S_b=\sum_{a\to b}t_a,\qquad
 w_b=\max\{r_b,m_b\}.
\]
Empty maxima are zero.  All fan slots of every selected top are
included in these definitions, whether or not a polygon of that top
visits the fan.  In particular
\begin{equation}\label{diagram:slot-sum}
 \sum_b S_b=7T.
\end{equation}
Write $J_b$ for the number of partial instances over $b$, and let
$P_b^*$ count the selected loose top sides belonging to turns at this
fan.  This number is zero unless $b$ has type $D$.  Each such loose side
records one incidence with the exceptional radial port $X_D$.

Let $I$ count internal contacts between two letter sides of the
fragment.  Radial side contacts are not counted in $I$.  Denote the
numbers of letter and radial fragment boundary sides by $B_{\rm raw}$
and $B_{\rm rad}$, and put $B=B_{\rm raw}+B_{\rm rad}$.

We first use the coincidence bound to control internal letter contacts.
The remaining estimate is local at fan apices: full links contribute
no error, while partial links contribute an error paid by radial
boundary sides.  Combining the two estimates will bound the fragment's
area by its total boundary length.

\begin{lemma}\label{diagram:weighted-rank}
Suppose Equation~\eqref{diagram:rank-bound} holds for collections of at
most $8A$ occurrences.  Then
\begin{equation}\label{diagram:contact-bound}
 I\le pT+f\sum_b(w_b-S_b)
              +\sum_e(h_e-\delta)t_e.
\end{equation}
\end{lemma}

\begin{proof}
Give each fan field of $b$ weight $w_b$, each loose field of $a$ weight
$t_a$, and each field of a decoration $e$ weight $t_e$.  These are
upper bounds for their side multiplicities in the fragment.  Indeed a
fan field is used at most once at a given normalized fan-apex link, so
its multiplicity is at most $r_b$.  A simple top polygon uses its loose
field at most once, and a copy of a decoration uses each indexed field
once.

For each integer $j\ge1$, apply
Equation~\eqref{diagram:rank-bound} to the positions of weight at least
$j$.  This is an eligible collection: $t_a\ge j$ implies
$w_b\ge m_b\ge t_a\ge j$ for every slot $a\to b$.  Its size is at most
$8A$, since there are at most $T$ distinct tops, at most $7T+Q$ fans,
and at most $E$ decorations.  Sum the resulting inequalities over $j$.
Within an unsigned-value class, the sum of its ranks is its total
weight minus its largest position weight.  Thus, writing
$\mathcal C$ for the unsigned-value classes, we obtain
\begin{equation}\label{diagram:weighted-coincidence}
 \sum_{C\in\mathcal C}
 \left(\sum_{v\in C}\operatorname{wt}(v)
                  -\max_{v\in C}\operatorname{wt}(v)\right)
 \le (p-7f)T+f\sum_b w_b+\sum_e(h_e-\delta)t_e.
\end{equation}

Every internal letter contact joins two different field positions.
For fan triangles and decorations this follows from the cancellations
in Lemma~\ref{diagram:normal-form}; for two loose top sides it follows
from the same-occurrence top cancellation.  A contact between different
kinds of cells automatically has different positions.  In each class
$C$, choose one maximum-weight position.  Charge each internal contact
to one of its endpoints other than that position.  Such an endpoint
always exists.  Different contacts use different actual side copies,
so the charge to any position is at most its side multiplicity and
hence at most its weight.  This proves that $I$ is bounded by the left
side of Equation~\eqref{diagram:weighted-coincidence}.  Using
Equation~\eqref{diagram:slot-sum} on its right side gives
Equation~\eqref{diagram:contact-bound}.
\end{proof}

This argument allows arbitrary cycles and repetitions in the contact
graph.  It charges side copies, and does not require the contacts to
form a forest.

\begin{lemma}\label{diagram:apex}
For every fan occurrence $b$ in a fragment,
\begin{equation}\label{diagram:apex-bound}
 \sum_{z\mapsto b}d_z+2(w_b-S_b)
 \le \frac12\left(\sum_{z\mapsto b}q_z+P_b^*\right)+2J_b.
\end{equation}
Moreover
\begin{equation}\label{diagram:partial-bound}
 2\sum_bJ_b\le B_{\rm rad}.
\end{equation}
\end{lemma}

\begin{proof}
First suppose $w_b=r_b$.  A simple top polygon has at most one turn at
this fan, so $S_b\ge\sum_{z\mapsto b}d_z$.  Consequently the left side
of Equation~\eqref{diagram:apex-bound} is at most
$\sum_{z\mapsto b}(2-d_z)$.
For a full link with $d_z=0$, at least four fan triangles give the
required inequality.  For a full link with $d_z=1$, the triangles form
a path in the fan between the two distinct ports of that top turn.
Ordinary ports are distinct line vertices, so this path has at least
two edges.  The only exception is a path incident to the port $X_D$;
it may have one edge, in which case its top turn has the loose side
counted by $P_b^*$.  Thus in either case its triangle count plus its
loose-side payment is at least two.  For $d_z\ge2$ the expression
$2-d_z$ is nonpositive.  At a represented partial link with $d_z=0$
there is at least one selected fan triangle, so its deficit is at most
$3/2$.  With $d_z=1$ its deficit is at most one, and with $d_z\ge2$
it is nonpositive.  Thus the error $2J_b$ suffices in this case.

Now suppose $w_b=m_b>r_b$, and choose a top occurrence $a_0$ with
$t_{a_0}=m_b$.  Call its turns dominant.  Let $x$ and $y$ be the total
numbers of selected dominant and nondominant turns over $b$.
Nondominant selected faces have at most one turn each at this fan,
including when other selected faces omit the fan entirely; hence
$y\le S_b-m_b$.  It follows that
\begin{equation}\label{diagram:dominant-algebra}
 \sum_{z\mapsto b}d_z+2(m_b-S_b)
       =x+y-2(S_b-m_b)\le x-y.
\end{equation}

Consider a full link.  In its cyclic list of turns the number of
consecutive dominant-to-dominant interfaces is at least the number of
dominant turns minus the number of other turns.  To see this, assign to
each dominant turn its following turn: at most as many dominant turns
can be followed by a nondominant one as there are nondominant turns.
A single turn is followed by itself.  Between a counted pair of turns
there is a path of fan triangles.  This path cannot have length zero,
because then the two top polygons use the same radial side of the
same occurrence $a_0$, forbidden by
Lemma~\ref{diagram:normal-form}.  If the path has length one, it must
have an exceptional port, since the other ports lie in the line part
of the bipartite fan graph.  That port has its loose-side payment.
Thus the path length plus the available payments is at least two for
each counted interface.  The triangle paths are disjoint portions of
the link.  A port incidence is an endpoint of just one such intervening
path, so no loose-side payment is used twice.  The triangle counts and
payments at full links therefore bound twice their total
dominant-minus-nondominant count from below.

At a partial link, the selected steps form a nonempty proper subset
of a cyclic list of at most five steps.  Split this subset into maximal
consecutive selected runs.  Each run and each intervening unselected
gap contains a step, so there are at most two runs.  In a run containing
turns, let $k$ and $l$ be the numbers of dominant and nondominant turns.
Their linear turn list has at most $l+1$ blocks of dominant turns.
Consequently, if $H$ counts its consecutive dominant-to-dominant
interfaces, then
\[
 H\ge k-l-1,\qquad k-l\le H+1.
\]
The fan-triangle path at each of these interfaces is entirely selected.
The same indexed-side cancellation and exceptional-port argument used
for a full link shows that its length plus its available loose-side
payments is at least two.  Different interfaces, including those in
different runs, use disjoint triangle paths and distinct port incidences,
so no payment is used twice.  A run without turns contributes zero to
the dominant-minus-nondominant count.  Summing over the at most two runs
therefore bounds that count by half the selected triangle count and
loose-side payments, plus two.

Choose the dominant occurrence $a_0$ once for the entire fan occurrence
$b$, and apply these estimates separately at all its apex instances.
They do not require the fragment to be connected, or every selected
polygon of $a_0$ to visit this fan.  Summing the partial- and full-link
estimates, and then using
Equation~\eqref{diagram:dominant-algebra}, proves
Equation~\eqref{diagram:apex-bound} also in the second case.

Finally, a represented partial cyclic link has at least two
transitions between selected and unselected steps.  Each transition
is a radial edge with exactly one selected incident face, and hence
is a radial fragment boundary side.  A radial side belongs to only one
fan-apex instance.  This proves
Equation~\eqref{diagram:partial-bound}.
\end{proof}

\begin{lemma}\label{diagram:fragment}
Put
\begin{equation}\label{diagram:kappa}
 \kappa=\min\{\beta-7\alpha,\alpha/5,\delta\}>0.
\end{equation}
If Equation~\eqref{diagram:rank-bound} holds for collections of at most
$8A$ occurrences, then every fragment of area $A$ satisfies
\begin{equation}\label{diagram:fragment-bound}
 A\le C_0 B,\qquad C_0=\frac{1}{\kappa}.
\end{equation}
\end{lemma}

\begin{proof}
Write $W=\sum_b w_b$ and $P^*=\sum_bP_b^*$.  Every simple top polygon
has at least three turns.  Summing
Equation~\eqref{diagram:apex-bound}, and using
Equation~\eqref{diagram:slot-sum}, gives
\begin{equation}\label{diagram:nominal-apex}
 3T+2(W-7T)\le \frac{Q+P^*}{2}+2\sum_bJ_b.
\end{equation}
The right side of Equation~\eqref{diagram:contact-bound} can be written
as
\[
 3T+2(W-7T)+\sum_eh_et_e-\Delta,
 \qquad
 \Delta=(\beta-7\alpha)T+\alpha W+\delta E.
\]
There are at most five selected triangles at a represented fan-apex
instance, so
$Q\le5\sum_b r_b\le5W$.  Therefore
\begin{equation}\label{diagram:defect}
 \Delta\ge\kappa(T+Q+E)=\kappa A.
\end{equation}
The total number of raw side copies in the fragment is
\[
 R_{\rm raw}=Q+P^*+\sum_e2h_et_e.
\]
Equations~\eqref{diagram:contact-bound} and
\eqref{diagram:nominal-apex} imply
\[
 I\le \frac{R_{\rm raw}}2-\Delta+2\sum_bJ_b.
\]
Since $B_{\rm raw}=R_{\rm raw}-2I$,
Equations~\eqref{diagram:partial-bound} and
\eqref{diagram:defect} give
\[
 2\kappa A\le 2\Delta
       \le B_{\rm raw}+4\sum_bJ_b
       \le B_{\rm raw}+2B_{\rm rad}\le2B.
\]
This is Equation~\eqref{diagram:fragment-bound}.
\end{proof}

\subsection{From bounded fragments to all pictures}

The probability estimate controls only bounded collections of
occurrences. To use it for an arbitrarily large picture, we partition
the faces into bounded sets while cutting few contacts. A fragment
was allowed to be disconnected and to have holes, so any partition
of the dual graph gives eligible fragments. The following separator
estimate makes their total new boundary small enough to absorb into
the area bound.

\begin{lemma}\label{diagram:separator}
For every $D<\infty$ and $\eta>0$ there is an integer $K=K(D,\eta)$
such that the vertices of any finite planar multigraph of maximum
degree at most $D$ can be partitioned into sets of size at most $K$,
with at most $\eta$ times the total number of vertices edges crossing
between sets.
\end{lemma}

\begin{proof}
Ignore loops and apply the planar separator theorem to the underlying
simple graph.  In a graph with $q$ vertices it gives a set of at most
$c\sqrt q$ vertices whose removal leaves two sets, with no edges
between them, each of size at most $2q/3$; one may take
$c=2\sqrt2$ \cite[Corollary~2]{LiptonTarjan1979}.  Repeat inside every set larger
than $K$, and make every removed separator vertex a singleton part.

Charge $c/\sqrt q$ to every vertex present in a recursive set of size
$q$.  These charges pay for all removed separator vertices.  For any
fixed original vertex, the successive set sizes containing it decrease
by a factor at most $2/3$ while they exceed $K$.  Reading that sequence
backward shows that its total charge is at most
\[
 \frac{c}{\sqrt K}\sum_{j\ge0}(\sqrt{2/3})^j
       =\frac{c}{(1-\sqrt{2/3})\sqrt K}.
\]
If the original graph has $A$ vertices, the number of removed
separator vertices is at most $C A/\sqrt K$, where
$C=c/(1-\sqrt{2/3})$.  Every edge crossing between final parts has a
removed separator vertex as an endpoint.  The degree bound therefore
limits their number, counting multiplicity, to $DC A/\sqrt K$.
Choose $K$ so large that $DC/\sqrt K\le\eta$.  Disconnected graphs
are treated component by component.
\end{proof}

\begin{proof}[Proof of Theorem~\ref{thm:diagram}]
Fix the parameters, so that $\kappa$ and $C_0$ in
Lemma~\ref{diagram:fragment} are positive constants.  Choose
$\eta>0$ with $2C_0\eta\le1/2$, and apply
Lemma~\ref{diagram:separator} with $D=15$.  This gives a fixed finite
threshold $K$.  By Lemma~\ref{diagram:rank}, with probability tending
to one Equation~\eqref{diagram:rank-bound} holds for all collections of
at most $8K$ occurrences.  Work with any realization having that
property.  The remainder of the proof is deterministic and applies
simultaneously to every picture for this realization.

For a normal picture, the dual graph of its interior faces and their
contacts is planar of degree at most fifteen.  Fan triangles have
three sides, decorations have at most six, and simple top polygons
have at most fifteen.  Parallel contacts and self-contacts cause no
problem in Lemma~\ref{diagram:separator}.  Partition its $A$ interior
faces into fragments of size at most $K$, cutting at most $\eta A$
contacts.  Let $B_0$ count contacts with the outside polygon; in the
sphere problem set $B_0=0$.  The sum of fragment boundary lengths is
at most $B_0+2\eta A$.  Lemma~\ref{diagram:fragment} applies to every
fragment, and summing it gives
\[
 A\le C_0(B_0+2\eta A),
 \qquad\text{hence}\qquad A\le2C_0 B_0.
\]
In a disk picture whose outside word has length $\ell$, one has
$B_0\le\ell$; outside-to-outside pairings represent free boundary
cancellations and do not add an interior contact.  Thus the normal
picture uses at most $2C_0\ell$ interior faces.

If a nonzero cellular two-cycle $c$ existed in $\widetilde P$, use its
normal spherical picture supplied by
Lemmas~\ref{diagram:existence} and \ref{diagram:normal-form}.
For it $B_0=0$, forcing $A=0$.  Its cellular chain would then be zero,
contrary to its construction.  Therefore
$H_2(\widetilde P;\mathbb Z)=0$.  The universal cover is a simply
connected two-dimensional CW complex, so all its reduced homology
groups vanish.  If it had a first nonzero homotopy group, the
Hurewicz theorem would identify it with a nonzero homology
group~\cite[Theorem~4.32]{Hatcher2002}.
It is consequently weakly contractible, and the Whitehead theorem
makes it contractible~\cite[Theorem~4.5]{Hatcher2002}.
This proves asphericity.

We finish by making explicit the claimed filling inequality over
letters.  Choose a vertex of every fan template and, for each of its
vertices, a fixed path from the chosen vertex to that vertex.  These
paths have length at most two.  Send every fan-apex vertex of a normal
picture to the bouquet vertex, and replace its indexed radial edge by
the label word of the corresponding fixed fan path, with the
appropriate orientation.  A fan triangle boundary becomes a closed
walk in that fan of length at most five.  A top polygon boundary
becomes a closed walk in its copy of $\Gamma$: at an ordinary point,
its two radial paths concatenate to a fan path between the two ports;
at the exceptional point the path to $X_D$ is followed by the loose
edge.  Its expanded length is at most thirty.  A decoration boundary
is unchanged.

Use as a finite letter presentation all simple cycle relators in each
included fan and top graph, together with the decoration rows.  This
presents the same group by the cone construction in
Section~\ref{sec:model}.  Any closed graph walk of length at most
thirty splits at repeated vertices into at most thirty simple cycles,
after free backtracks are erased.  Each expanded face above thus has
a filling with at most thirty relator disks in this presentation.
Replacing the faces of the normal picture by these fillings leaves
the outside letter word unchanged and produces a disk diagram, with
free folds allowed, of area at most
$60C_0\ell$.  Sphere components can be discarded.  Any other fixed
finite choice of defining cycle relators gives the same linear
inequality after bounded relator substitutions.  This proves the
theorem.
\end{proof}

\begin{remark}\label{diagram:quantifiers}
The separator threshold is chosen after the positive parameters are
fixed and before $n$ tends to infinity.  No uniform bound as
$\alpha,\beta,\delta\longrightarrow0$ is asserted or required.
For each fixed parameter choice, the single finite-collection event in
Lemma~\ref{diagram:rank} controls all disk words and all cellular
two-cycles simultaneously.
\end{remark}
